\subsubsection{Características do RCBD}

O experimento foi estruturado como um delineamento em blocos casualizados com os seguintes parâmetros:

\begin{align*}
\text{Número de Blocos} &= 3 \\
\text{Número de Tratamentos} &= 4 \\
\text{Total de Observações} &= 60
\end{align*}

\paragraph{Modelo RCBD}
O modelo estatístico é dado por:

\[Y_{ij} = \mu + \tau_i + \beta_j + \epsilon_{ij}\]

onde: $Y_{ij}$ é a observação no tratamento $i$ e bloco $j$; $\mu$ é a média geral; $\tau_i$ é o efeito do tratamento $i$; $\beta_j$ é o efeito do bloco $j$; $\epsilon_{ij}$ é o erro aleatório.

\paragraph{Benefícios do Bloqueamento}
O uso de blocos reduz a variância residual, isolando a variação sistemática devida aos blocos do erro aleatório. Isto aumenta o poder estatístico para detectar diferenças entre tratamentos.

\paragraph{Interpretação dos Resultados}
Espera-se que:

\begin{enumerate}
\item O efeito do bloco seja significativo (p-valor baixo)
\item O quadrado médio residual (MSResíduo) seja menor que em um delineamento completamente casualizado
\item O p-valor do tratamento seja menor, indicando maior poder para detectar diferenças
\end{enumerate}

